\documentclass[a4paper]{article}
\usepackage{amsfonts}
\usepackage{amsmath}
\usepackage{amssymb}
\usepackage{graphicx}     

\begin{document}
  
\noindent \large Auteur: Sadia Nazary  \\
\noindent \large Studierichting: Informatica\\
\noindent \large Oefeningenreeks 1D nr. 22.7\\

\medskip

\normalsize

\textbf{Opgave:} Bepaal de nulpunten van de veelterm in $\mathbb{C}$ en ontbind in factoren over $\mathbb{C}[z]$: 
\[
z^4 - 2z^3 - 2z^2 + 8z - 8
\]

\textbf{Oplossing:}\\

We proberen een paar eenvoudige waarden in te vullen:

\[
z = 2 \Rightarrow 2^4 - 2 \cdot 2^3 - 2 \cdot 2^2 + 8 \cdot 2 - 8 = 16 - 16 - 8 + 16 - 8 = 0
\]

Dus $z = 2$ is een nulpunt. We delen de veelterm door $(z - 2)$:\\

\begin{tabular}{c|ccccc}
       & $1$ & $-2$ & $-2$ & $8$ & $-8$ \\
$2$    &     & $2$  & $0$  & $-4$ & $8$ \\ \hline
       & $1$ & $0$  & $-2$ & $4$ & $0$
\end{tabular}

\[
\Rightarrow z^4 - 2z^3 - 2z^2 + 8z - 8 = (z - 2)(z^3 - 2z + 4)
\]

We testen opnieuw:\\

\[
z = -2 \Rightarrow (-2)^3 - 2(-2) + 4 = -8 + 4 + 4 = 0
\]

Dus $z = -2$ is ook een nulpunt. We delen verder:\\

\begin{tabular}{c|cccc}
       & $1$ & $0$ & $-2$ & $4$ \\
$-2$   &     & $-2$ & $4$ & $-4$ \\ \hline
       & $1$ & $-2$ & $2$ & $0$
\end{tabular}

\[
\Rightarrow z^3 - 2z + 4 = (z + 2)(z^2 - 2z + 2)
\]

We lossen $z^2 - 2z + 2 = 0$ op:\\

\[
z = \frac{2 \pm \sqrt{(-2)^2 - 4 \cdot 1 \cdot 2}}{2} = \frac{2 \pm \sqrt{-4}}{2} = 1 \pm i
\]

\textbf{Eindresultaat:}\\

\[
z^4 - 2z^3 - 2z^2 + 8z - 8 = (z - 2)(z + 2)(z - (1 + i))(z - (1 - i))
\]

\textbf{Nulpunten:} \\

\[
z_0 = 2, \quad z_1 = -2, \quad z_2 = 1 + i, \quad z_3 = 1 - i
\]

\begin{thebibliography}{999}
%\bibitem{a} Referentie
\end{thebibliography}
\end{document}
